\clearpage
\begin{examplebox}
\textbf{例1｜只求范围：焦点弦同支与异支最短分别是多少？}

对于 $H:x^2/9-y^2/16=1$，有 $a=3,b=4,c=5$。
\begin{center}% 精确坐标：H: x^2/9-y^2/16=1, c=5, b^2/a=16/3.
% 两图中蓝色曲线由解析式 x=+-3 sqrt(1+y^2/16) 绘制，非手绘。
\begin{tikzpicture}[x=0.39cm,y=0.39cm,>=stealth, font=\small]
\begin{scope}
  \draw[->,gray] (-6.8,0)--(7.3,0) node[right]{$x$};
  \draw[->,gray] (0,-7.8)--(0,7.8) node[above]{$y$};
  \draw[thick,blue!75!black,domain=-7:7,samples=90] plot ({3*sqrt(1+\x*\x/16)},{\x});
  \draw[thick,blue!75!black,domain=-7:7,samples=90] plot ({-3*sqrt(1+\x*\x/16)},{\x});
  \draw[very thick,red!75!black] (5,-5.33333333)--(5,5.33333333);
  \fill (5,0) circle (2.1pt) node[below right]{$F_2$};
  \fill[red!75!black] (5,5.33333333) circle (1.8pt) node[above left=5pt]{$A$};
  \fill[red!75!black] (5,-5.33333333) circle (1.8pt) node[below left=5pt]{$B$};
  \node[align=center] at (0,-9.0) {同支焦点弦\quad $x=c$\\最短 $2b^2/a$};
\end{scope}
\begin{scope}[xshift=6.6cm]
  \draw[->,gray] (-6.8,0)--(7.3,0) node[right]{$x$};
  \draw[->,gray] (0,-7.8)--(0,7.8) node[above]{$y$};
  \draw[thick,blue!75!black,domain=-7:7,samples=90] plot ({3*sqrt(1+\x*\x/16)},{\x});
  \draw[thick,blue!75!black,domain=-7:7,samples=90] plot ({-3*sqrt(1+\x*\x/16)},{\x});
  \draw[very thick,red!75!black] (-3,0)--(3,0);
  \fill (5,0) circle (2.1pt) node[above right]{$F_2$};
  \fill[red!75!black] (-3,0) circle (1.8pt) node[above left=5pt]{$A$};
  \fill[red!75!black] (3,0) circle (1.8pt) node[above right=5pt]{$B$};
  \node[align=center] at (0,-9.0) {异支焦点弦\quad $y=0$\\最短 $2a$};
\end{scope}
\end{tikzpicture}
\end{center}
\begin{itemize}
\item 异支：最短 $2a=6$，在 $y=0$ 取到，交点 $(-3,0),(3,0)$；
\item 同支：最短 $2b^2/a=32/3$，在 $x=5$ 取到，交点 $(5,\pm16/3)$；
\item 两条平行渐近线方向的排除斜率是 $k=\pm4/3$。
\end{itemize}
\textbf{识别信号：}题目只问最值，不必联立方程；\textbf{先看同支或异支}。
\end{examplebox}

\clearpage
\begin{examplebox}
\textbf{例2｜代入斜率：过右焦点的直线斜率为 $2$，求焦点弦长。}

仍取 $x^2/9-y^2/16=1$，$k=2$，有 $|k|>4/3$，是\textbf{右支同支弦}。
\[
L=\frac{2\cdot3\cdot16(1+2^2)}{|9\cdot2^2-16|}
=\frac{480}{20}=\boxed{24}.
\]
\textbf{检查：}$24>32/3$，和同支下界一致。
\textbf{扩展：}$k=1/2$ 时为异支弦，
\[
L=\frac{96(1+1/4)}{16-9/4}=\boxed{\frac{480}{55}=\frac{96}{11}}>6.
\]
\end{examplebox}

\clearpage
\begin{examplebox}
\textbf{例3｜逆向求斜率：异支焦点弦长为 $8$，求斜率。}

还是 $a=3,b=4$，$0\le u=k^2<16/9$。
\[
8=\frac{96(1+u)}{16-9u}
\Longrightarrow 128-72u=96+96u
\Longrightarrow u=\frac4{21}.
\]
所以 $\boxed{k=\pm\dfrac2{\sqrt{21}}}$。
\textbf{有效性核验：}$u=4/21<16/9$，确为异支；
若忘记先判断类型，就容易把分母绝对值误拆成另一支。
\end{examplebox}

\clearpage
\begin{examplebox}
\textbf{例4｜高价值变式：长度给定时，同支/异支可能同时有解吗？}

对 $x^2/9-y^2/16=1$，假设 $L=12$。
异支公式 $12=96(1+u)/(16-9u)$ 导出
\[
192-108u=96+96u\Longrightarrow \boxed{u=\frac8{17}}<\frac{16}{9}.
\]
同支公式 $12=96(1+u)/(9u-16)$ 导出
\[
108u-192=96+96u\Longrightarrow \boxed{u=24}>\frac{16}{9}.
\]
故既有异支直线 $k=\pm\sqrt{8/17}$，也有同支直线 $k=\pm2\sqrt6$。
\textbf{解释：}同一长度 $L$ 可能对应两类焦点弦，需按 $2a$ 和 $2b^2/a$ 两个下界分别判断。

\textbf{一题多解检查：}给定长度 $L$ 时，可用
\[
\begin{aligned}
\boxed{u_{\text{异支}}=\frac{b^2(L-2a)}{a^2L+2ab^2}},&\quad L\ge2a;\\[4pt]
\boxed{u_{\text{同支}}=\frac{b^2(L+2a)}{a^2L-2ab^2}},&\quad L>\frac{2b^2}{a}.
\end{aligned}
\]
同支等号 $L=2b^2/a$ 时只有\textbf{竖直线}，不能用第二式求有限斜率。
\end{examplebox}
